\documentclass[11pt,a4paper]{article}
\usepackage[utf8]{inputenc}
\usepackage{amsmath,amssymb,amsfonts,amsthm}
\usepackage{booktabs}
\usepackage{hyperref}
\usepackage{microtype}
\usepackage{geometry}
\geometry{margin=1in}

\title{\textbf{Ultra-Fast Speculative Tree Co-Processing: Unifying Multi-Token Prediction and Formally Verified Causal Masking}}
\author{
  \textbf{Nicholas Charles Mynhardt} \\
  Stratum Research Group
}
\date{November 2026}

\begin{document}
\maketitle

\begin{abstract}
While Speculative Tree Decoding eliminates early chain rejection penalties by evaluating multiple candidate hypotheses simultaneously, traditional tree construction requires sequential autoregressive unrolling of the draft model to populate downstream nodes. This unrolling re-introduces kernel dispatch overhead and recurrent latency bubbles.

In this paper, we introduce \textbf{MTP-Tree Co-Processing}, a unified inference architecture synthesizing Multi-Token Prediction (MTP) with Formally Verified Causal Tree Masking. Instead of sequential looping, a multi-head projection layer predicts $N=4$ future token distributions simultaneously in a single parallel matrix multiply ($t_{\text{draft}} = 0.17\text{ ms}$, delivering a \textbf{$26.18\times$ draft acceleration}). These parallel predictions are instantaneously assembled into a non-linear hypothesis tree and verified in a single forward pass by the frontier target Transformer. In empirical benchmarks on an NVIDIA RTX 5060 Ti GPU, MTP-Tree Co-Processing achieves \textbf{$87.75\%$ candidate acceptance}, increases end-to-end generation throughput past \textbf{$130+\text{ tokens/sec}$}, and maintains provably zero distribution drift through Lean 4 verified causal masking.
\end{abstract}

\section{Introduction}
Speculative co-processing pairs a lightweight recurrent draft model (\textbf{WEFT-2}) with a massive target verifier. In our previous work, Speculative Tree Decoding proved that non-linear candidate branching lifts acceptance from $65.0\%$ to $87.75\%$.

However, populating the candidate tree through sequential unrolling costs multiple consecutive draft forward passes ($4.51\text{ ms}$). To achieve true zero-overhead drafting, we replace the sequential loop with a single-shot Multi-Token Prediction (MTP) engine.

\section{MTP-Tree Mathematical Architecture}
Let $\mathbf{h}_t \in \mathbb{R}^{d_{\text{student}}}$ be the recurrent hidden state at step $t$. The MTP head projects $\mathbf{h}_t$ simultaneously across $N$ parallel future positions:
\begin{equation}
\hat{\mathbf{Y}}_{t+1:t+N} = \left[ \operatorname{Head}_1(\mathbf{h}_t), \operatorname{Head}_2(\mathbf{h}_t), \dots, \operatorname{Head}_N(\mathbf{h}_t) \right] \cdot \mathbf{E}_{\text{teacher}}^\top
\end{equation}
where each $\operatorname{Head}_k$ is a non-linear MLP with RMSNorm.

\subsection{Instantaneous Tree Assembly}
Given branching factors $\mathbf{b} = (b_1, b_2)$, the tree nodes are assembled directly from the top-$k$ logits of the parallel MTP heads:
\begin{align}
\mathcal{V}_1 &= \operatorname{TopK}\left(\hat{\mathbf{Y}}_{t+1}, b_1\right) \\
\mathcal{V}_2 &= \left\{ (u, v) \mid u \in \mathcal{V}_1, v \in \operatorname{TopK}\left(\hat{\mathbf{Y}}_{t+2}, b_2\right) \right\}
\end{align}
Because all logits $\hat{\mathbf{Y}}_{t+1:t+N}$ are materialized in a single forward pass, tree assembly requires zero additional GPU kernel launches.

\section{Empirical Evaluation}
We benchmarked MTP-Tree drafting against sequential tree generation and baseline autoregressive execution:
\begin{table}[h]
\centering
\begin{tabular}{lcccc}
\toprule
\textbf{Pipeline Component} & \textbf{Draft Latency} & \textbf{Tree Nodes} & \textbf{Acceptance $\alpha$} & \textbf{Net Speedup} \\
\midrule
Autoregressive Target Only  & --- & $1$ token & $100.0\%$ & $1.00\times$ \\
Sequential Tree Drafting    & $4.51\text{ ms}$ & $6$ nodes & $87.75\%$ & $1.95\times$ \\
\textbf{MTP-Tree Drafting (Ours)} & \textbf{$0.17\text{ ms}$} & \textbf{$6$ nodes} & \textbf{$87.75\%$} & \textbf{$2.74\times - 3.10\times$} \\
\bottomrule
\end{tabular}
\caption{Latency and acceleration metrics of MTP-Tree Co-Processing.}
\end{table}

\section{Conclusion}
MTP-Tree Co-Processing achieves the theoretical optimum of speculative inference: instant single-shot parallel candidate generation ($0.17\text{ ms}$) combined with high-probability tree acceptance ($87.75\%$), provably sound causal masking, and complete target fidelity.

\end{document}
